% Assembles the repetitions of the previous extended abstract; does not replace the proofs.
\subsection*{Map of results and dependencies}
The argument contains five levels that should be distinguished when reading. The arithmetic construction supplies the modes; their readouts define the complete form; the estimates establish positivity on precise domains; the Weil criterion determines which additional identification permits RH to be concluded. The order of presentation makes the principal result visible, while the body preserves the order of its dependencies.

\begin{center}
\renewcommand{\arraystretch}{1.22}
\begin{tabular}{@{}p{0.21\linewidth}p{0.48\linewidth}p{0.23\linewidth}@{}}
\hline
\textbf{Level} & \textbf{Result and role in the argument} & \textbf{Location}\\
\hline
Common construction &
Histories, sheets, memory, and incidences remain within a single refinement system. The center, lifted boundary, and vacancies specify the information transmitted by each readout. &
Sections~\ref{lectura:manuscrito-05} and~\ref{sec:centro-frontera}.\\
Arithmetic and functional &
Native product, normal forms, and irreducibles; occupations, Euler product, prime--power moments, finite parts, and theta--Mellin continuation. &
Sections~\ref{lectura:manuscrito-01}--\ref{lectura:manuscrito-04};
beginning of~\ref{lectura:manuscrito-03}.\\
Reconstruction and transport &
The columns jointly preserve the prime, gamma, and polar terms. The gamma tail reconstructs the test function. Compression, memory, and Cayley preserve their balances and cross products. &
Sections~\ref{lectura:manuscrito-03},
\ref{sec:transporte-momentos}, and~\ref{sec:lectores-gamma}.\\
Sign of the form &
Coercivity of the specified assembly, control of details, and Schur reduction. The complete window of length \(1/9\) satisfies
\(\mathscr W(f,f)\geq\frac12\|f\|_2^2\). &
Section~\ref{lectura:manuscrito-03} and
Theorems~\ref{schur-add-schur-exacto},
\ref{schur-add:coercividad-ventana}.\\
Global criterion &
Identification of the complete moments, with all their cross terms and coverage of the test domain, would permit application of the Weil criterion. It is not replaced by the preceding local bounds. &
Section~\ref{lectura:manuscrito-03};
subsections~\ref{tm-add-c5} and~\ref{tm-add-m7}.\\
\hline
\end{tabular}
\end{center}

Two distinctions accompany this map. In arithmetic, the period of a residual action does not replace irreducibility of the product. In analysis, a positive factorization of the memory moments is not automatically identified with that of the arithmetic moments. The constants \(8/81\) and \(5/9\) belong, respectively, to the compression balance and the modal memory ratio; their proofs preserve these distinct roles.

The center \((5,5)\), the regional signature \(555555\), its three continuation sheets, and the schedules \(21/22\) and \(6/7\) are explained in their constructive positions. The readouts \(72,27,77,22\), the lifted boundary, and the completion \(729+271\) are not grouped by numerical coincidence: each preserves the operation and domain that produces it. Their role in this article is to specify the transportable information before passing to the spectral problem.

Composition of readouts also preserves two distinct operations: the
diagonal product of residual tables and the tensor product of their
moments, which produces divisor convolution. Stieltjes transport includes
the affine contribution of the pole. The Bendersky--Glaisher hierarchy
relates regularized derivatives at negative even levels to positive odd
values of the same functional. Enclosures from irreducible-mode cutoffs
complement the depth bounds, with their respective errors explicit
(Sections~\ref{ix-add:corte-primos},
\ref{ix-add:composicion-lectores} and~\ref{ix-add:bendersky}).

The definitions and proofs used are brought together in the indicated sections and Appendices A--C, with the domain of each result stated explicitly. The local proofs, functional results and conditional statements retain their own scope here.
